Excessive groundwater withdrawal is a principal cause of land subsidence in many aquifer-dependent regions worldwide \citep{galloway_land_1999,herrera-garcia_mapping_2021,jasechko_rapid_2024,huning_global_2024}. The resulting loss of land elevation can damage infrastructure, increase exposure to inundation, and reduce agricultural productivity \citep{tzampoglou_selected_2023}. By 2040, as much as 19\% of the global population may live in areas with a high probability of subsidence, with much of this exposure concentrated in coastal plains and alluvial deltas \citep{herrera-garcia_mapping_2021,nicholls_global_2021,buffardi_issue_2023}. Prominent cases associated with groundwater withdrawal include the Central Valley, Mexico City, and Jakarta, where cumulative displacement has reached several meters \citep{miller_rapid_2020,chaussard_over_2021,chaussard_sinking_2013}. The Mekong Delta presents a related regional hazard, with widespread subsidence occurring at rates of centimeters per year \citep{erban_groundwater_2014}.

Effective mitigation depends on observations that reveal both land-surface displacement and deformation within the aquifer system. Spirit leveling provides accurate surface measurements but is commonly limited to annual or semiannual surveys because of the field work involved \citep{poland_guidebook_1984}. Continuous Global Navigation Satellite System stations record three-dimensional displacement at fixed locations, while Interferometric Synthetic Aperture Radar provides repeated measurements over broad areas \citep{bock_physical_2016,galloway_application_2007}. These methods differ in temporal and spatial coverage, but each describes deformation at the land surface. They cannot determine how that deformation is distributed among subsurface depth intervals \citep{zhang_inverse_2016,burbey_extensometer_2020}. Observations of compaction at depth are therefore needed to connect surface displacement with deformation within the aquifer system.

Multilayer compaction monitoring wells (MLCWs) provide this depth-resolved information. Magnetic rings anchored at different depths record relative displacement within a single borehole, allowing deformation to be separated among monitored aquifer and aquitard intervals \citep{hung_measuring_2021}. The resulting records reveal where compaction is concentrated within the subsurface profile rather than only its combined expression at the land surface. When interpreted together with hydraulic head, elastic and inelastic deformation measured by an MLCW can also support estimates of skeletal specific storage and safe groundwater levels \citep{hung_measuring_2021}. MLCWs therefore add direct evidence of deformation at depth that surface geodetic observations cannot provide.

An MLCW requires a dedicated borehole, investigations of subsurface materials, and repeated observations throughout its operation \citep{hung_measuring_2021,hung2025_realtime}. Field personnel must collect the measurements, and the records must be processed before the latest depth-resolved deformation can be interpreted. Access to current MLCW information may therefore be delayed, separated by longer measurement intervals, or eventually interrupted. Hydraulic head and cGNSS observations may remain available each month during these intervals. Hydraulic head reflects changes in pore-water pressure, whereas cGNSS records the integrated vertical response at the surface. Although neither record identifies where deformation occurs within the subsurface profile, their continued availability may support monthly estimates between direct MLCW observations. Establishing the limits of these estimates could inform how monitoring resources are allocated while preserving MLCW observations as the direct measure of deformation with depth.

Surface displacement and hydraulic observations have been combined to infer how deformation is distributed with depth and to estimate aquifer-system properties \citep{jiang2018_insar,smith2021_depth,verberne2025_ravenna,verberne2025_dutch}. In central Taiwan, MLCW records have also supported analyses of aquifer-system compaction and groundwater management \citep{hung2012_mlcw,hung_measuring_2021}. More recent data-driven methods have reconstructed missing compaction records or examined changes under reduced groundwater use \citep{liu_reconstructing_2023,liu_deep_2025}. A related study used a high-frequency extensometer record to forecast short-term vertical deformation and compared extensometer and MLCW observations at corresponding depths \citep{hung2025_realtime}. This body of work demonstrates the value of combining deformation and hydrologic records, but it does not resolve how monthly estimates for MLCW-defined depth intervals behave when new MLCW information is delayed, collected less frequently, or no longer available. The extent to which monthly hydraulic head and cGNSS observations can support such estimates therefore remains uncertain.

The present study addressed this uncertainty using records from one MLCW station within a regional monitoring network in central Taiwan. Separate relations were estimated for six standardized depth sections using changes in hydraulic head and vertical surface displacement from the same month and preceding months, together with seasonal terms derived from the calendar date. Because several of these quantities may rise or fall together, their separate relations to deformation cannot always be determined precisely. Bayesian ridge regression reduces the weight given to relations that are not well supported by the observed record and provides a range of plausible values for each monthly estimate. Historical MLCW observations were used to establish these relations and assess the estimates, but previous MLCW deformation increments were not included among the model inputs. Borehole lithology was examined separately to interpret material differences among the depth sections. The assessment considered both estimation error and the uncertainty surrounding each estimate.

The primary evaluation examined whether monthly hydraulic-head and cGNSS observations could support depth-resolved deformation estimates while finalized MLCW records remained temporarily unavailable. It also examined how estimation error varied among depth sections and over the interval between successive data releases. The analysis then considered less frequent MLCW measurements after initial records of different lengths, focusing on monthly error, accumulated discrepancy, and uncertainty. A final evaluation treated the absence of subsequent MLCW observations as a limiting case and traced the monthly and accumulated discrepancies over time.

Section 2 describes the monitoring setting and datasets used in the analysis, and Section 3 presents the analytical approach and experimental designs. Section 4 reports the results, while Section 5 examines their implications for monthly depth-resolved deformation monitoring. Section 6 summarizes the principal conclusions and identifies the questions that remain open.
